% Generated draft: requires mathematical and source review.
\documentclass[11pt]{article}
\usepackage[margin=25mm]{geometry}
\usepackage{fontspec}
\setmainfont{Latin Modern Roman}
\usepackage{amsmath,amssymb,mathtools,unicode-math}
\setmathfont{Latin Modern Math}
\usepackage{enumitem,needspace,xurl,hyperref}
\setlength{\parindent}{0pt}
\setlength{\parskip}{5pt}
\setlength{\emergencystretch}{3em}
\newcommand{\N}{\mathbb{N}}
\newcommand{\Z}{\mathbb{Z}}
\newcommand{\Q}{\mathbb{Q}}
\newcommand{\R}{\mathbb{R}}
\newcommand{\C}{\mathbb{C}}
\newcommand{\F}{\mathbb{F}}
\newcommand{\A}{\mathbb{A}}
\newcommand{\PP}{\mathbb{P}}
\newcommand{\E}{\mathbb{E}}
\newcommand{\eps}{\varepsilon}
\newcommand{\dd}{\,\mathrm{d}}
\begin{document}
\section*{UnsolvedMath ID: 700205}
\textbf{Source problem code:} catalogue entry
\subsection{Definitions and mathematical statement}
A group \(G\) is soluble if its derived series, defined by \(G^{(0)}=G\) and \(G^{(k+1)}=[G^{(k)},G^{(k)}]\), reaches the trivial group after finitely many steps. It is non-soluble otherwise. A group is simple if it is nontrivial and its only normal subgroups are the trivial subgroup and itself. These conventional group definitions do not determine the meaning of the numerical classes named in the supplied source.

\textbf{Problem as supplied.} Determine the laws of distribution of non-soluble and simple group-theoretic numbers in the sequence of natural numbers. The source additionally refers to Problem 2.78.

\textbf{Unresolved qualifications.} The supplied record neither defines a ``group-theoretic number'' nor explains when such a number is called non-soluble or simple. In particular, it does not specify whether membership concerns the existence of a group of a given order, a condition on every group of that order, or another construction. The referenced Problem 2.78 is not supplied. Furthermore, ``laws of distribution'' specifies neither a particular statistic nor a proposed asymptotic statement. Thus the numerical classes and the precise distributional target cannot be defined unambiguously from this record. The convention concerning whether zero belongs to the natural numbers is also unstated.
\subsection{Short English statement}
P. I. Trofimov asks how non-soluble and simple group-theoretic numbers are distributed among the natural numbers. The defining terminology and the referenced Problem 2.78 are missing, so a precise numerical formulation remains unresolved.
\subsection{Sources}
\begin{itemize}
\item Local UnsolvedMath registry: supplied record with permanent ID \(700205\), exact source code \texttt{catalogue entry}, titled \emph{Kourovka Notebook Problem 3.57}. The registry labels the problem open; this is supplied metadata, not an independently verified status assessment.
\item Supplied bibliographic attribution: \emph{The Kourovka Notebook (Unsolved Problems in Group Theory)}, Problem 3.57, third issue (1969); proposed by P. I. Trofimov. The record identifies the twenty-first issue edition (2026), \texttt{arXiv:1401.0300v48}. Its collection description names E. I. Khukhro and V. D. Mazurov as editors and covers Issues 1--21 (1965--2026).
\item The statement refers to Problem 2.78, whose text is absent from the supplied record. The cited source and cross-reference have not been independently consulted.
\end{itemize}
\end{document}
